\section*{Capítulo \ref{ch:b1:titration-methods} --- Métodos y curvas de valoración}

\begin{solution}{exo:b1:titration-methods:1}
\ce{H3O+ + OH- -> 2H2O}, $K = 1/K_e = 10^{14.00}$. $C = 12.4 \times
0.100/20.0 = \qty{0.0620}{mol/L}$.
\end{solution}

\begin{solution}{exo:b1:titration-methods:2}
En la \omterm{def:b1:titration-methods:titration}{equivalencia}, la disolución es cloruro de amonio de
\qty{0.050}{mol/L}: $\mathrm{pH} = \frac12(9.25 - \log 0.050) = 5.28$. El
rojo de metilo (3.8--5.8) lo contiene; el naranja de metilo, cuya zona
(2.5--4.5) solo se alcanza después de la \omterm{def:b1:titration-methods:titration}{equivalencia}, al bajar el pH,
cambiaría demasiado tarde.
\end{solution}

\begin{solution}{exo:b1:titration-methods:3}
\ce{MnO4- + 5Fe^2+ + 8H+ -> Mn^2+ + 5Fe^3+ + 4H2O};
$C = 5 \times 0.0200 \times 12.6/10.0 = \qty{0.126}{mol/L}$.
\end{solution}

\begin{solution}{exo:b1:titration-methods:4}
$C = 0.0100 \times 14.2/50.0 = \qty{2.84e-3}{mol/L}$ de calcio más
magnesio. Ambos reaccionan antes del \omterm{def:b1:titration-methods:titration}{punto final}, que solo da su suma:
una \omterm{def:b1:titration-methods:kinds}{valoración simultánea}.
\end{solution}

\begin{solution}{exo:b1:titration-methods:5}
0~mL: \omterm{def:b1:predominant-reaction:strong-acid}{ácido débil}, $\mathrm{pH} = \frac12(4.76 + 1.00) = 2.88$. 5.0~mL:
semiequivalencia, 4.76. 10.0~mL: 8.73
(\cref{prop:b1:titration-methods:ph-at-equivalence}). 15.0~mL: exceso de
hidróxido de \qty{0.50}{mmol} en \qty{25.0}{mL}, $[\ce{OH-}] = 0.020$,
$\mathrm{pH} = 12.30$. Todos concuerdan con la curva calculada con un
margen de 0.01.
\end{solution}

\begin{solution}{exo:b1:titration-methods:6}
Los $\mathrm{p}K_a$ 2.15 y 7.21 difieren en más de 4, igual que 7.21 y
12.34: dos \omterm{def:b1:titration-methods:titration}{equivalencias}, a 10.0 y a \qty{20.0}{mL}. Primera: disolución
de \ce{H2PO4-}, $\mathrm{pH} \approx \frac12(2.15 + 7.21) = 4.68$;
segunda: \ce{HPO4^2-}, $\frac12(7.21 + 12.34) = 9.78$ (la curva exacta da
4.71 y 9.66, pues las fórmulas desprecian la dilución y el agua). La
tercera acidez ($\mathrm{p}K_a$ 12.34) está demasiado cerca de la del
agua: el hidróxido añadido solo reacciona en parte y el pH sube sin
salto.
\end{solution}

\begin{solution}{exo:b1:titration-methods:7}
Ácido clorhídrico: $6.2 \times 0.100/10.0 = \qty{0.062}{mol/L}$; ácido
etanoico: $(14.6 - 6.2) \times 0.100/10.0 = \qty{0.084}{mol/L}$. A medio
camino, la mitad del ácido etanoico está valorada:
$\mathrm{pH} = \mathrm{p}K_a = 4.76$.
\end{solution}

\begin{solution}{exo:b1:titration-methods:8}
$V_{\mathrm{eq}} = \qty{10.0}{mL}$. 2.0~mL: $E = 0.77 + 0.059\log(2/8) =
0.73$~V; 9.0~mL: $0.77 + 0.059\log 9 = 0.83$~V; 11.0~mL: exceso de
permanganato de \qty{0.020}{mmol}, \ce{Mn^2+} \qty{0.200}{mmol}, $E = 1.51 +
\frac{0.059}{5}\log 0.10 = 1.50$~V; 20.0~mL: 1.51~V.
\end{solution}

\begin{solution}{exo:b1:titration-methods:9}
0~mL: $\kappa = (349.8 + 76.2) \times 0.0100 = \qty{4.26}{mS/cm}$.
10.0~mL: \ce{Na+} y \ce{Cl-} a $1.00/110 = \qty{9.09e-3}{mol/L}$,
$\kappa = (50.3 + 76.2) \times \num{9.09e-3} = 1.15$. 20.0~mL: \ce{Na+}
$0.0167$, \ce{Cl-} y \ce{OH-} $0.00833$~mol/L, $\kappa = 0.84 + 0.635 + 1.65
= 3.13$~mS/cm. Pendientes (despreciando la dilución, \qty{1}{mL} aporta
\qty{1.0e-3}{mol/L}): $(50.3 - 349.8) \times 10^{-3} = -0.30$ y
$(50.3 + 198.3) \times 10^{-3} = +0.25$~mS/cm por mL.
\end{solution}

\begin{solution}{exo:b1:titration-methods:10}
En $V_{\mathrm{eq}} - v$ el exceso de ácido es $C_Bv$, y en
$V_{\mathrm{eq}} + v$ el exceso de base es $C_Bv$; en el mismo volumen,
$[\ce{H3O+}]$ antes es igual a $[\ce{OH-}]$ después, de modo que
$\mathrm{pH}(V_{\mathrm{eq}} - v) +
\mathrm{pH}(V_{\mathrm{eq}} + v) = 14$: simetría respecto de
(V$_{\mathrm{eq}}$, 7). Salto: \qty{1.0e-5}{mol} de exceso en
\qty{20}{mL}, pH de 3.30 a 10.70, 7.4 unidades. Divididas por 100:
$\num{1.0e-7}$~mol en \qty{20}{mL}, pH de 5.3 a 8.7, 3.4 unidades: un
\omterm{def:b1:titration-methods:titration}{punto final} mucho menos nítido.
\end{solution}

\begin{solution}{exo:b1:titration-methods:11}
Antes: \ce{Fe^3+} formado $= 5C_{\mathrm{Mn}}V$, \ce{Fe^2+} restante $=
5C_{\mathrm{Mn}}(V_{\mathrm{eq}} - V)$, de modo que $E = 0.77 + 0.059\log\frac{V}{V_{\mathrm{eq}}
- V}$: en $V_{\mathrm{eq}}/2$, $E = 0.77$. Después: exceso de \ce{MnO4-}
$\propto
V - V_{\mathrm{eq}}$, \ce{Mn^2+} $\propto V_{\mathrm{eq}}$: $E = 1.51 +
\frac{0.059}{5}\log\frac{V - V_{\mathrm{eq}}}{V_{\mathrm{eq}}}$ (pH 0); en
$2V_{\mathrm{eq}}$, 1.51~V. En la \omterm{def:b1:titration-methods:titration}{equivalencia}, $(0.77 + 5 \times 1.51)/6 =
1.39$~V. A pH 1, el par del permanganato baja $0.059 \times 8/5 =
0.094$~V: $E_{\mathrm{eq}} = (0.77 + 5 \times 1.416)/6 = 1.31$~V.
\end{solution}

\begin{solution}{exo:b1:titration-methods:12}
\ce{NH4+ + OH- -> NH3 + H2O} tiene $K = 10^{14.00 - 9.25} = 10^{4.75}$:
la reacción no es completa cerca de la \omterm{def:b1:titration-methods:titration}{equivalencia} y el salto es
pequeño, en torno a pH 11 (pH de \omterm{def:b1:titration-methods:titration}{equivalencia}
$14 - \frac12(4.75 + 1.30) = 11.0$). \omterm{def:b1:titration-methods:kinds}{Valoración por retroceso}: se añade
$n_1$ de hidróxido de sodio (en exceso), se hierve para eliminar el
amoniaco formado y después se valora el hidróxido restante con ácido
clorhídrico ($n_2$, salto nítido fuerte--fuerte):
$n(\ce{NH4+}) = n_1 - n_2$.
\end{solution}

\begin{solution}{pb:b1:titration-methods:1}
\textbf{1.} \ce{C6H6O6 + 2H+ + 2e- -> C6H8O6}.
\textbf{2.} \ce{C6H8O6 + I2 -> C6H6O6 + 2I- + 2H+}.
\textbf{3.} 0 y $-$I.
\textbf{4.} \ce{I2 + 2S2O3^2- -> 2I- + S4O6^2-}, $\log K = 2(0.53 - 0.02)/0.059
= 17$.
\textbf{5.} El color azul necesita yodo; desaparece con el último yodo,
en la \omterm{def:b1:titration-methods:titration}{equivalencia}.
\textbf{6.} Las pérdidas por el aire rebajarían el resultado; añadido de
una vez y en exceso, el yodo oxida el ácido ascórbico de forma rápida y
completa.
\textbf{7.} Una reacción rápida y \omterm{def:b1:extent-q-and-k:quantitative}{cuantitativa} durante toda la
\omterm{def:b1:titration-methods:titration}{valoración}, sin oxidación por el aire mientras se añade lentamente el
\omterm{def:b1:titration-methods:titration}{valorante}, y un \omterm{def:b1:titration-methods:titration}{punto final} con el primer exceso de yodo.
\textbf{8.} Exceso de yodo $n_R$; reacción con el ácido ascórbico;
\omterm{def:b1:titration-methods:titration}{valoración} del yodo restante con tiosulfato.
\textbf{9.} De lo contrario, quedaría ácido ascórbico, y el exceso de
yodo, nulo, no diría cuánto. Aquí reaccionaron \qty{2.775e-4}{mol} de
\qty{5.000e-4}{mol}: efectivamente, en exceso.
\textbf{10.} Las disoluciones de tiosulfato cambian lentamente con el
tiempo; su concentración interviene directamente en el resultado.
\textbf{11.} 10.
\textbf{12.} Una pipeta aforada de \qty{20.00}{mL}.
\textbf{13.} $20.00 \times 10^{-3} \times 0.0250 = \qty{5.000e-4}{mol}$.
\textbf{14.} $8.90 \times 10^{-3} \times 0.0500 = \qty{4.450e-4}{mol}$.
\textbf{15.} La mitad: \qty{2.225e-4}{mol}.
\textbf{16.} $5.000 - 2.225 = \qty{2.775e-4}{mol}$.
\textbf{17.} Mol a mol: \qty{2.775e-4}{mol}.
\textbf{18.} $\times 10$: \qty{2.775e-3}{mol}, $\times 176.0$:
\qty{0.488}{g}.
\textbf{19.} Alrededor de un \qty{2}{\percent} por debajo de la etiqueta.
\textbf{20.} $n = \bigl(C_IV_I - \frac12C_TV_T\bigr)\,V_{\text{matraz}}/V_{\text{alícuota}}$.
\textbf{21.} Yodo introducido: pipeta $0.03/20.00 = \qty{0.15}{\percent}$
y concentración \qty{0.2}{\percent}, alrededor de un \qty{0.25}{\percent},
es decir, \qty{1.3e-6}{mol}. Exceso: bureta $0.05/8.90 = \qty{0.56}{\percent}$
y concentración \qty{0.2}{\percent}, alrededor de un
\qty{0.60}{\percent}, es decir, \qty{1.3e-6}{mol}.
\textbf{22.} Las incertidumbres absolutas se combinan, mientras que la
diferencia es menor que cada una de las cantidades:
$\sqrt{1.3^2 + 1.3^2} \times 10^{-6} =
\qty{1.8e-6}{mol}$, un \qty{0.66}{\percent} de \qty{2.775e-4}{mol}, más
que cualquiera de los dos términos.
\textbf{23.} Con el matraz (\qty{0.1}{\percent}) y la pipeta de la
alícuota
(\qty{0.2}{\percent}): $\sqrt{0.66^2 + 0.1^2 + 0.2^2} = \qty{0.70}{\percent}$.
\textbf{24.} $0.488 \pm 0.003$~g: \textbf{\qty{0.49}{g} de ácido ascórbico
por comprimido}.
\end{solution}
